\scanpage{613}\mspage{1}
% archive sheet number: 612

1. \uline{Rappels sur les couples de Quillen}. \uline{Couples de Quillen spéciaux}.

Je vais reprendre les notations et la terminologie pour les relations de
LLP et RLP de Quillen. J'écris (pour une catégorie $M$, et
$\Phi,\Psi\subset\mathrm{Fl}(M)$)
\[
 \Phi \mathrel{\tikz[baseline=-.45ex,x=1em,y=1ex]{\draw (0,0)--(2.15,0) (0.32,-0.5)--(0,0) (1.83,0.5)--(2.15,0);}} \Psi \Longleftrightarrow (\Phi\subset\Psi^*)
 \Longleftrightarrow(\Psi\subset\Phi_*) \tag*{(1.1)}
\]
i.e. $\Phi$ a le LLP p.r. à $\Psi$, ou encore $\Psi$ le \deleted{\unsure{R}}RLP
p.r. à $\Phi$. Je \deleted{écris} dis aussi que $(\Phi,\Psi)$ est un
\uline{couple} Q-\uline{orthogonal}. Cette notion n'est pas symétrique,
mais je me permets parfois (exceptionellement) d'écrire
\[
 \Psi \mathrel{\tikz[baseline=-.45ex,x=1em,y=1ex]{\draw (0,0)--(2.15,0) (0.32,0.5)--(0,0) (1.83,-0.5)--(2.15,0);}} \Phi \qquad\text{comme synonyme de}\qquad
 \unsure{\text{``}}\Phi \mathrel{\tikz[baseline=-.45ex,x=1em,y=1ex]{\draw (0,0)--(2.15,0) (0.32,-0.5)--(0,0) (1.83,0.5)--(2.15,0);}} \Psi .
 \tag*{(1.2)}
\]
$\Phi$ est Q-\uline{orthogonal à gauche} à $\Psi$, $\Psi$
\uline{est Q-orthogonal à droite} à $\Phi$, et la Q-orthogonalité
gauche \inserted{(resp. à droite)} est notée par le
\inserted{demi-}sommet des flèches pointant vers le bas
(resp. vers le haut). J'écris
\[
\left\{\begin{aligned}
\Phi \mathrel{\tikz[baseline=-.45ex,x=1em,y=1ex]{\draw (0,0)--(3,0) (0.32,0.5)--(0,0) (0.32,-0.5)--(0,0) (2.68,0.5)--(3,0);}} \Psi &\Longleftrightarrow
 (\Phi\unsure{=|\deleted{=}}\unsure{\Psi^*|\deleted{\Psi^*}}\deleted{\unsure{\Phi|\Psi}}\unsure{)|\deleted{)}}\\
\Phi \mathrel{\tikz[baseline=-.45ex,x=1em,y=1ex]{\draw (0,0)--(3,0) (0.32,-0.5)--(0,0) (2.68,0.5)--(3,0) (2.68,-0.5)--(3,0);}} \Psi &\Longleftrightarrow (\Psi=\Phi_*)\\
\Phi \mathrel{\tikz[baseline=-.45ex,x=1em,y=1ex]{\draw (0,0)--(3,0) (0.32,0.5)--(0,0) (0.32,-0.5)--(0,0) (2.68,0.5)--(3,0) (2.68,-0.5)--(3,0);}} \Psi &\Longleftrightarrow
 (\Phi=\Psi^*\text{ et }\Psi=\Phi_*).
\end{aligned}\right.\tag*{(1.3.)}
\]
Ces trois cas de Q-orthogonalité renforcée sont donc ceux où on a
de plus $\Phi=\Phi^{\sim}$ resp. $\Psi=\overline\Psi$ resp.
$\Phi=\Phi^{\sim}$ et $\Psi=\overline\Psi$. \deleted{\unsure{O}}On dit que $\Phi$ est
Q-\uline{clos à gauche} si $\Phi=\Phi^{\sim}$, que $\Psi$ est
\uline{Q-clos à droite} si $\Psi=\Psi^{-}$, que le couple Q-orthogonal
$(\Phi,\Psi)$ est Q-\uline{clos à gauche} (resp. \uline{à droite})
si $\Phi$ \deleted{\unsure{ou}} (resp. $\Psi$) est Q-clos à g.
(resp. à droite), on dit qu'il est Q-clos si les deux conditions
